\documentclass[11pt]{article}

\usepackage[margin=1in]{geometry}
\usepackage{amsmath,amssymb,amsthm,mathtools}
\usepackage{booktabs}
\usepackage{tabularx}
\usepackage{enumitem}
\usepackage{xcolor}
\usepackage{hyperref}
\usepackage[T1]{fontenc}

\hypersetup{colorlinks=true,linkcolor=blue,citecolor=blue,urlcolor=blue}

\newtheorem{theorem}{Theorem}[section]
\newtheorem{proposition}[theorem]{Proposition}
\newtheorem{lemma}[theorem]{Lemma}
\newtheorem{corollary}[theorem]{Corollary}
\theoremstyle{definition}
\newtheorem{definition}[theorem]{Definition}
\newtheorem{remark}[theorem]{Remark}

\newcommand{\Z}{\mathbb Z}
\newcommand{\F}{\mathcal F}
\newcommand{\kg}{\mathrm{KG}}
\newcommand{\MSF}{\mathrm{MSF}}
\DeclareMathOperator{\LCM}{lcm}

\title{Erd\H{o}s Problem \#7: Odd Distinct Covering Systems\\
\large BFF reconstruction, finite certificates, and prime-power refinements}
\author{Research working draft}
\date{14 August 2026}

\begin{document}
\maketitle

\begin{center}
\fcolorbox{red!60!black}{red!5}{\parbox{0.92\textwidth}{
\textbf{Status warning.} This is a working research document, not a paper or novelty claim.
Statements explicitly labelled ``published'' are traced to cited sources.  Statements labelled
``derived'' are deductions made in this project and should be independently checked before
circulation.  The finite computations are reproduced by the accompanying exact-rational Python
scripts.  In particular, the bound $218{,}243{,}025$ and the maximum-spanning-forest refinement
have not been peer reviewed and are not asserted here to be absent from the prior literature.
}}
\end{center}

\begin{abstract}
Erd\H{o}s and Selfridge asked whether there exists a finite covering system of congruences
$a_i\pmod{m_i}$ whose moduli are pairwise distinct, greater than one, and all odd.  We organize
the problem in finite cyclic-group, product-space, divisor-lattice, hypergraph, probabilistic,
and optimization languages.  We then reconstruct the exponent-sensitive necessary inequality
of Berger, Felzenbaum and Fraenkel (BFF), emphasizing that their correction term is a forest
bound on forced intersections.  This immediately suggests a derived maximum-spanning-forest
refinement.  Combining published structural restrictions with exact enumeration of exponent
vectors gives a provisional, fully reproducible deduction that a hypothetical odd covering has
LCM at least $218{,}243{,}025$.  The boundary shape
$3^3 5^2 7\cdot11\cdot13\cdot17\cdot19$ survives both the published BFF inequality and the
first forest/triangle refinements, identifying a concrete target for further work.  A block or
``tower capacity'' lemma is introduced as a second route for exploiting prime-power depth.  We
then condition directly on the $3$-adic coordinate while retaining the actual cofactor residue
classes.  This replaces the earlier reciprocal-capacity diagnostic by CRT-aware set bounds,
forces a common residual of at least $234523/3233230$ in one nine-fibre branch, and yields a
second $675\times323323$ ``heavy-fibre'' necessary condition for the boundary shape.  A further
conditioning on the $7$-coordinate writes the boundary shape as $4725\times46189$, produces a
global common residual in the square-free four-prime cofactor, and reduces exclusion of the
boundary case to a finite residual-concentration inequality on only $4725$ outer points.
\end{abstract}

\tableofcontents

\section{The problem and finite reformulations}

A \emph{covering system} is a finite family
\[
 a_i\pmod{m_i},\qquad 1\le i\le k,
\]
whose union is all of $\Z$.  It is \emph{distinct} when the $m_i>1$ are pairwise distinct.
Erd\H{o}s Problem \#7 asks whether a distinct covering can have every $m_i$ odd.

Let
\[
 L=\LCM(m_1,\dots,m_k).
\]
The covering is periodic modulo $L$, so it is equivalent to a cover of $\Z/L\Z$ by cosets of
subgroups.  If
\[
 L=\prod_{i=1}^{n}p_i^{s_i},\qquad p_1<\cdots<p_n,
\]
then the Chinese remainder theorem gives
\[
 \Z/L\Z\cong\prod_{i=1}^{n}\Z/p_i^{s_i}\Z.
\]
A divisor $m=\prod p_i^{d_i}$ is naturally encoded by a depth vector
$\mathbf d=(d_1,\dots,d_n)$, $0\le d_i\le s_i$.  Thus a residue class modulo $m$ is a cylinder
fixing $d_i$ $p_i$-adic digits in coordinate $i$.  The square-free case is the special case
$s_i=1$ for every $i$; Balister et al. proved that an odd distinct covering is impossible in
that case \cite{BalisterSquarefree2021}.  Consequently the unresolved problem is intrinsically
multilevel or prime-power in nature.

The same finite object can be viewed as a constrained set-cover problem, a hypergraph cover,
or a probability problem on uniformly random $X\in\Z/L\Z$.  For the event
$A_i=\{X\equiv a_i\pmod{m_i}\}$, $\Pr(A_i)=1/m_i$, while the Chinese remainder theorem controls
many intersections.  These translations motivate both the BFF overlap inequality and modern
SAT/MIP approaches.

\section{Published structural restrictions}

We record only the restrictions used later.

\begin{theorem}[BFF II, published \cite{BFF1987}]
If an odd incongruent covering system exists and the LCM has $n$ distinct prime factors, then
$n\ge6$.  More precisely, the exponent-sensitive parameters defined in Section~\ref{sec:bff}
must satisfy $g(\bar w,\bar{\mathbf z})\ge2$.
\end{theorem}

\begin{theorem}[Hough--Nielsen, published \cite{HoughNielsen2019}]
Every distinct covering system contains a modulus divisible by $2$ or $3$.
Hence, for an odd distinct covering,
\[
 3\mid L.
\]
\end{theorem}

\begin{theorem}[Balister et al., published \cite{BalisterSquarefree2021}]
There is no distinct covering system with all moduli odd and square-free.  In particular, the
LCM of a hypothetical odd distinct covering cannot be square-free.
\end{theorem}

A different theorem of Balister et al. gives the useful restriction that an odd covering LCM
must be divisible by $9$ or $15$ \cite{BalisterDensity2022}.  It is not required for the
finite bound derived below, but it is consistent with the first surviving shape.

\section{Re-deriving the BFF inequality}\label{sec:bff}

\subsection{The product box and building blocks}

Write
\[
 L=\prod_{i=1}^{n}p_i^{s_i},\qquad P_i=p_i^{s_i}.
\]
BFF represent the relevant finite group by the product box
\[
 \mathcal P=[0,P_1)\times\cdots\times[0,P_n),
 \qquad |\mathcal P|=L,
\]
and refine cosets into product ``building blocks'' indexed by nonempty subsets
$I\subseteq[n]$.  The index records coordinates on which the block is restricted.
After enlarging the family in a way that can only make covering easier, the number of
blocks of a given index is bounded by BFF's quantity $\alpha(I)$.  The exceptional counts
for one-coordinate indices are the source of the asymmetric first coordinate in the final
formula \cite{BFF1987}.

Remove the blocks whose index has cardinality one.  The remaining set $\mathcal S$ is again a
product set.  Its coordinate sizes are
\[
 y_1=P_1-\frac{P_1-1}{p_1-1}
     =\frac{(p_1-2)P_1+1}{p_1-1},
\]
and, for $i\ge2$,
\[
 y_i=P_i-2\frac{P_i-1}{p_i-1}
     =\frac{(p_i-3)P_i+2}{p_i-1}.
\]
This normalization makes the BFF parameters transparent:
\begin{align}
 w&=\frac{P_1-1}{(p_1-2)P_1+1},\\
 z_1&=\frac{p_1^{s_1-1}-1}{(p_1-2)P_1+1},\\
 z_i&=\frac{P_i-1}{(p_i-3)P_i+2},\qquad 2\le i\le n.
\end{align}
They are normalized capacities of higher-index blocks inside the surviving box.

\subsection{The pre-forest union bound}

For $|I|\ge2$, BFF bound the fraction of $\mathcal S$ occupied by all blocks of index $I$.
Summing these bounds and collecting terms gives
\begin{equation}\label{eq:g1}
 g_1(w,\mathbf z)
 =(1+w)\prod_{i=2}^{n}(1+z_i)-w
 -(1+w-z_1)\sum_{i=2}^{n}z_i.
\end{equation}
The naive argument would require $g_1\ge2$: the higher-index blocks must cover the entire
residual product set, and $g_1-1$ is their normalized union-bound capacity.

\subsection{Forced intersections and the forest lemma}

Let $\mathcal J=\binom{[n]}2$.  For $I,J\in\mathcal J$ with $I\cap J=\varnothing$, BFF show
under their worst-case enlargement that the corresponding type-$I$ and type-$J$ unions have
forced intersection of normalized size
\begin{equation}\label{eq:beta}
 \beta(I,J)=\prod_{r\in I\cup J}z_r.
\end{equation}
They use the elementary forest inequality
\[
 \left|\bigcup_{v\in V}S_v\right|
 \le \sum_{v\in V}|S_v|-
 \sum_{\{u,v\}\in E(T)}|S_u\cap S_v|
\]
for any forest $T$ on the set family.  Therefore every forest on $\mathcal J$ whose edges join
disjoint 2-subsets gives a legitimate forced-overlap correction.

BFF choose one particular forest supported on the first five coordinates, with total weight
\begin{equation}\label{eq:wbff}
 W_{\rm BFF}
 =3z_1z_2z_3z_4+3z_1z_2z_3z_5
  +2z_1z_2z_4z_5+z_1z_3z_4z_5.
\end{equation}
Thus their published polynomial is
\begin{equation}\label{eq:gbff}
 g_{\rm BFF}=g_1-W_{\rm BFF},
\end{equation}
and a necessary condition for an odd distinct cover is
\begin{equation}\label{eq:bffcondition}
 g_{\rm BFF}\ge2.
\end{equation}
This re-derivation highlights the conceptual content: BFF is a naive capacity estimate minus a
rigorously selected set of unavoidable overlaps.

\section{Derived maximum-spanning-forest refinement}

The published forest lemma does not require BFF's particular forest.  This yields an immediate
optimization.

\begin{definition}
Let $G_n$ be the Kneser disjointness graph $\kg(n,2)$: vertices are the 2-subsets of $[n]$ and
$I$ is adjacent to $J$ precisely when $I\cap J=\varnothing$.  Give each edge the weight
$\beta(I,J)$ from \eqref{eq:beta}.  Let $W_{\MSF}(\mathbf z)$ be the maximum total weight of a
forest in $G_n$.
\end{definition}

\begin{proposition}[Derived]
Every odd distinct covering satisfies
\begin{equation}\label{eq:gmsf}
 g_{\MSF}(w,\mathbf z)
 :=g_1(w,\mathbf z)-W_{\MSF}(\mathbf z)\ge2.
\end{equation}
Moreover $g_{\MSF}\le g_{\rm BFF}$.
\end{proposition}

\begin{proof}
BFF's forest lemma and their exact intersection formula apply to every forest with edges between
disjoint 2-index sets.  Choose a maximum-weight such forest.  The BFF forest is one feasible
choice, hence $W_{\MSF}\ge W_{\rm BFF}$ and the resulting necessary condition is no weaker.
\end{proof}

For $n\ge5$, $\kg(n,2)$ is connected, so with positive weights the optimum can be taken as a
maximum-weight spanning tree.  The accompanying script `msf\_refinement.py` constructs it by
exact-rational Kruskal optimization.

As a minimal hand-check that BFF's displayed tree need not be optimal, one may replace the last
term of \eqref{eq:wbff} by $z_2z_3z_4z_5$ using a different admissible tree.  The gain is
$(z_2-z_1)z_3z_4z_5$ whenever $z_2>z_1$; otherwise retain the published tree.

\subsection{Connection with Hunter--Worsley bounds: where the pairwise route stops}

The forest inequality used by BFF is the same structural mechanism as the Hunter--Worsley
second-order upper bound for a union: subtract a maximum-weight spanning tree of pairwise
intersections from the sum of the singleton probabilities \cite{Hunter1976,Worsley1982}.
Modern probability-optimization literature sharpens the interpretation.  For fixed singleton
probabilities together with only lower bounds on pairwise intersections, the optimal upper union
bound is of Hunter type; Ramachandra and Natarajan summarize this result and its relation to
Boros et al. \cite{BorosEtAl2014,RamachandraNatarajan2023}.

Our BFF setting is not literally the generic fixed-marginal probability problem---the singleton
quantities in \eqref{eq:g1} arise from an arithmetic relaxation---so we do not import an
``optimality theorem'' blindly.  Nevertheless, after that relaxation, the only additional data
used by the forest step are lower bounds on pair intersections for disjoint 2-index groups.  The
maximum-spanning-forest refinement therefore exhausts the obvious Hunter--Worsley improvement of
that pairwise data.  This is an important negative conclusion: a substantial next improvement
should use information that BFF's pairwise forest discards, such as exact triple intersections,
compatibility among coordinate projections, higher-index groups, or the prime-power fibre
structure.

\section{A provisional exact lower bound for the LCM}\label{sec:218m}

This section combines published theorems with a finite exact computation.  It is deliberately
labelled provisional pending independent verification and a more exhaustive novelty search.

\subsection{Prime monotonicity for fixed exponent pattern}

For $i\ge2$ define
\[
 z(p,s)=\frac{p^s-1}{(p-3)p^s+2}.
\]
Treating $p$ temporarily as a real variable,
\begin{equation}\label{eq:zderivative}
 \frac{\partial z}{\partial p}
 =\frac{p^{s-1}\left(s(p-1)+p-p^{s+1}\right)}{\left((p-3)p^s+2\right)^2}.
\end{equation}
For $p\ge5$ and $s\ge1$, the numerator is negative, so $z(p,s)$ decreases with $p$.
BFF state that $g_{\rm BFF}$ is increasing in its variables on the relevant domain
\cite{BFF1987}.  Since Hough--Nielsen forces $p_1=3$, replacing any later ordered prime by a
smaller allowed prime can only increase $g_{\rm BFF}$ for a fixed exponent vector.  The
$s_1=1$ boundary, for which $z_1=0$, follows by evaluating the same polynomial on the closure of
the domain.

\subsection{Finite exponent enumeration}

Set
\begin{equation}\label{eq:X}
 X=3^3 5^2 7\cdot11\cdot13\cdot17\cdot19
  =218{,}243{,}025.
\end{equation}

For exactly six distinct prime factors, it suffices by the preceding monotonicity to enumerate
positive exponent vectors for the minimal tuple
\[
 (3,5,7,11,13,17)
\]
whose minimal-prime shape is below $X$.  There are 89 such vectors.  Exact rational evaluation
of \eqref{eq:gbff} gives $g_{\rm BFF}<2$ in every case.  The largest value is
\[
 \frac{506755}{255717}\approx1.98170242886
\]
at exponent vector $(4,3,1,1,1,1)$, whose minimal-prime shape is $172{,}297{,}125$.

For exactly seven distinct prime factors, the minimal tuple is
\[
 (3,5,7,11,13,17,19).
\]
There are 16 positive exponent vectors below $X$.  Again every one has $g_{\rm BFF}<2$; the
maximum is
\[
 \frac{117043}{60775}\approx1.92584121761
\]
at $(2,2,1,1,1,1,1)$, with shape $72{,}747{,}675$.

For eight or more distinct primes, nonsquarefreeness and $3\mid L$ imply
\[
 L\ge3^2\cdot5\cdot7\cdot11\cdot13\cdot17\cdot19\cdot23
   =334{,}639{,}305>X.
\]

\begin{proposition}[Derived + exact finite certificate; provisional]
Assuming the published BFF, Hough--Nielsen, and square-free theorems as quoted above, the
accompanying exact finite certificate yields
\[
 \boxed{L\ge218{,}243{,}025}
\]
for the LCM of any hypothetical odd distinct covering system.
\end{proposition}

\begin{proof}
BFF gives at least six distinct prime factors; Hough--Nielsen gives $3\mid L$; and the
square-free theorem forces at least one exponent to exceed one.  If $L<X$, then $n\ge8$ is
impossible by the displayed numerical lower bound.  For $n=6$ and $n=7$, prime monotonicity
reduces every fixed exponent pattern to the corresponding minimal prime tuple, and the complete
positive-exponent enumerations below $X$ have respectively 89 and 16 cases.  The accompanying
CSV records the exact fraction $g_{\rm BFF}$ for every case; all are strictly below two,
contradicting the BFF necessary condition.
\end{proof}

The script uses Python's `fractions.Fraction`; floating-point numbers are printed only for human
readability.  The CSV is therefore an explicit finite arithmetic certificate, though not a
formal proof-assistant certificate.

\section{The first surviving boundary shape}

At the boundary
\begin{equation}\label{eq:L0}
 L_0=3^3 5^2 7\cdot11\cdot13\cdot17\cdot19
     =218{,}243{,}025,
\end{equation}
the BFF parameters are
\[
 w=\frac{13}{14},\qquad
 (z_1,\dots,z_7)=
 \left(\frac27,\frac6{13},\frac15,\frac19,\frac1{11},\frac1{15},\frac1{17}\right).
\]
Exact evaluation gives
\[
 g_{\rm BFF}=\frac{175089}{85085}\approx2.05781277546>2,
\]
so the published inequality no longer excludes the shape.

The maximum-spanning-forest refinement gives
\[
 g_{\MSF}=\frac{2608283}{1276275}\approx2.04366848837>2.
\]
Thus forest optimization closes about a quarter of the BFF excess above two, but not enough to
eliminate $L_0$.

\section{Going one order beyond forests: an exact triangle cluster}

When $n\ge6$, $\kg(n,2)$ contains triangles: three pairwise-disjoint 2-subsets
$I,J,K$.  For such a triple, product structure fixes not only all three pair intersections but
also the triple intersection.  Hence the union of these three type groups admits exact
third-order inclusion--exclusion,
\[
 |A_I\cup A_J\cup A_K|
 =\sum |A|-\sum |A\cap A'|+|A_I\cap A_J\cap A_K|,
\]
with normalized triple intersection
\[
 \gamma(I,J,K)=\prod_{r\in I\cup J\cup K}z_r.
\]
Contract the triangle to one cluster.  The remaining groups can again be attached by a forest;
for an edge from the cluster to an outside group, use any one triangle member whose guaranteed
intersection lies inside the cluster intersection, and choose the largest such lower bound.
This gives a rigorous derived refinement.

At $L_0$, exhaustive exact search over triangle choices gives
\[
 g_{\triangle}=\frac{23472541}{11486475}
 \approx2.04349384820>2.
\]
The gain beyond the maximum spanning forest is real but small.  This experiment indicates that
higher-order type-2 inclusion--exclusion by itself is unlikely to close the remaining gap; the
prime-power depth should enter more directly.

\section{A block or tower capacity lemma}

McNew and Setty define $r(n)$ as the largest number of residues modulo $n$ coverable by distinct
nontrivial divisors of $n$, and $c(n)=1+r(n)/n$.  Thus $n$ is a covering number exactly when
$c(n)=2$.  They prove
\[
 c(mn)\le c(m)h(n),\qquad (m,n)=1,
 \qquad h(n)=\frac{\sigma(n)}n,
\]
and conjecture the stronger coprime submultiplicativity $c(mn)\le c(m)c(n)$
\cite{McNewSetty2026}.  The following elementary block lemma is tailored to the CRT product.

\begin{lemma}[Derived block-capacity certificate]\label{lem:block}
Let
\[
 N=M_1\cdots M_t,
\]
where the $M_j$ are pairwise coprime.  Suppose $R_j$ is an upper bound for $r(M_j)$.  Call a
nontrivial divisor $d\mid N$ \emph{pure} if $d\mid M_j$ for some $j$, and \emph{mixed}
otherwise.  If
\begin{equation}\label{eq:blockcertificate}
 \sum_{\substack{d\mid N\\d\text{ mixed}}}\frac1d
 <
 \prod_{j=1}^{t}\left(1-\frac{R_j}{M_j}\right),
\end{equation}
then $N$ is not a covering number.
\end{lemma}

\begin{proof}
Pure congruences supported inside block $M_j$ cover at most $R_j$ residues in that coordinate.
By CRT, after all pure congruences there remains a Cartesian product of coordinate complements
of relative size at least the right-hand side of \eqref{eq:blockcertificate}.  A mixed modulus
$d$ covers exactly a fraction $1/d$ of the full space, so all mixed moduli together cover at
most the left-hand side by the union bound.  If that capacity is smaller than the guaranteed
pure-block residual, some point remains uncovered.
\end{proof}

The left side is easy to compute:
\[
 \sum_{d\text{ mixed}}\frac1d
 =h(N)-1-\sum_{j=1}^{t}\bigl(h(M_j)-1\bigr).
\]
For prime powers one has the exact formula
\[
 r(p^a)=1+p+\cdots+p^{a-1}=\frac{p^a-1}{p-1},
\]
since the union bound is attained by choosing pairwise-disjoint residue classes at successive
$p$-adic depths.

At $L_0$, use the tower blocks
\[
 27,25,7,11,13,17,19.
\]
The guaranteed pure-block residual is
\[
 \frac{3072}{12155}\approx0.25273549979,
\]
whereas the mixed reciprocal budget is
\[
 \frac{6954221}{11486475}\approx0.60542690425.
\]
Thus the simplest block certificate fails.  This failure is informative: a successful
prime-power argument must either absorb many mixed moduli into larger blocks, improve the
block capacities $r(M)$ substantially, or prove significant unavoidable overlap among the
mixed moduli.

\section{Conditioning on the 3-adic coordinate with actual cofactor sets}\label{sec:cofactor}

The scalar calculation in the previous version of this draft kept only reciprocal capacities in
each fibre.  It showed why a naive fibre union bound cannot work, but it discarded the main
arithmetic information: in every compatible fibre a class modulo $3^a d$ projects to the
\emph{same residue class modulo $d$} in the cofactor coordinate.  We now retain those sets.

Write
\[
 L_0=27Q,\qquad Q=25\cdot7\cdot11\cdot13\cdot17\cdot19=8{,}083{,}075.
\]
For $t\in\Z/27\Z$, identify the fibre
\[
 F_t=\{x\in\Z/L_0\Z:x\equiv t\pmod{27}\}
\]
with $\Z/Q\Z$ by CRT.  If a congruence has modulus $3^a d$, $(d,3)=1$, then it either misses
$F_t$ or, on every compatible fibre, becomes one residue class
\[
 C_{a,d}\subseteq \Z/Q\Z
\]
of modulus $d$.  In particular, a depth-zero class $d$ gives the same set in all 27 fibres; a
depth-one class $3d$ gives the same set in all nine fibres below one mod-$3$ branch; and so on.
This is the set-valued coupling that the reciprocal-capacity model lost.

\subsection{CRT-aware bounds for the cofactor unions}

McNew and Setty define $r(n)$ as the largest number of residues modulo $n$ coverable by
distinct divisor moduli, and use a CRT inclusion--exclusion bound over pairwise-coprime
subfamilies \cite{McNewSetty2026}.  Their same calculation applies to the small-multiplicity
multisets needed below.  In the present cases the multiplicity is smaller than the least prime
divisor of the cofactor, so repeated copies of one modulus can be taken with distinct residues.
For $k$ copies of each nontrivial divisor, the normalized bound is
\[
 U_k(N)=\sum_{\substack{d\mid N\\d>1}}
 \frac{B(k,\omega(d))}{d},
 \qquad
 B(k,r)=-\sum_{j=1}^{r}(-k)^j S_2(r,j),
\]
where $S_2(r,j)$ is a Stirling number of the second kind.  The script
\texttt{cofactor\_set\_attack.py} evaluates these quantities exactly.

For $Q$ we obtain
\[
 U_1(Q)=\frac{3338}{5225}\approx0.63885167464,
\]
so the depth-zero cofactor classes leave an \emph{actual set}
\[
 R_0=(\Z/Q\Z)\setminus U_0
\]
with
\[
 \frac{|R_0|}{Q}\ge
 1-U_1(Q)=\frac{1887}{5225}\approx0.36114832536.
\]
This is much stronger than the old reciprocal-sum estimate, which allowed the depth-zero
classes nominal capacity $0.85561$.

If one ignores the 3-adic branch assignment and optimistically allows two copies of every
cofactor divisor in a single fibre, the same CRT calculation gives
\[
 U_2(Q)=\frac{1536162}{1616615}\approx0.95023366726.
\]
Hence every fibre not covered wholesale by the modulus-$3$ class still has a nonempty
cofactor residual after depths zero and one, of relative size at least
\[
 \frac{80453}{1616615}\approx0.04976633274.
\]
The branch coupling lets us improve this further.

\subsection{A common residual in one nine-fibre branch}

The unique modulus-$3$ congruence covers one complete mod-$3$ branch.  Any class $3d$ with
$d>1$ assigned to that same branch is globally redundant: moving its mod-$3$ residue to one
of the other two branches cannot uncover anything in the already-full branch and can only add
coverage elsewhere.  Thus, when deriving a necessary condition, we may optimistically assume
that every nontrivial depth-one class is assigned to one of the two remaining branches, call
them $B$ and $C$.

Let $U_B,U_C\subseteq\Z/Q\Z$ be the actual unions obtained from all depth-zero classes together
with the depth-one classes compatible with $B$ or $C$, respectively.  The same pairwise-coprime
CRT polynomial can be summed over the two branches before the residue choices are known.  If
$S$ is a pairwise-coprime collection of $k$ cofactor divisors and $x$ of their second copies are
sent to $B$, its multiplicity factor in the two branch bounds is
\[
 2^x+2^{k-x}.
\]
For odd $k$ (positive inclusion--exclusion sign), this is at most $2^k+1$.  For even $k$
(negative sign), it is at least $2^{k/2+1}$.  Applying those inequalities termwise gives the
assignment-independent bound
\[
 \frac{|U_B|+|U_C|}{Q}
 \le
 \frac{2998707}{1616615}
 \approx1.85492959053.
\]
Therefore one of the two branches, say $B$, has a \emph{common residual set}
\[
 R_B=(\Z/Q\Z)\setminus U_B
\]
shared by all nine of its descendant mod-$27$ fibres, with
\begin{equation}\label{eq:branch-residual}
 \boxed{
 \frac{|R_B|}{Q}\ge
 \frac{234523}{3233230}
 \approx0.07253520473.}
\end{equation}
This is the first genuinely set-coupled $3$-adic obstruction in the project: the same residual
points of the cofactor must now be covered nine times by the depth-two and depth-three classes.

\subsection{A second conditioning: $675\times R$ and heavy fibres}

The branch residual is not yet large enough to defeat the remaining classes by a raw capacity
argument.  A useful second view is obtained by grouping both prime-power coordinates:
\[
 L_0=675R,\qquad
 R=7\cdot11\cdot13\cdot17\cdot19=323{,}323.
\]
The $R$-coordinate is square-free.  First consider the ``pure base'' classes whose moduli are
nontrivial divisors of $675$ and have $R$-cofactor $1$.  The one-copy CRT bound gives
\[
 U_1(675)=\frac{487}{675}.
\]
Consequently at least
\[
 675-487=188
\]
of the $675$ base fibres are not covered wholesale by a pure $3^a5^b$ congruence.

Fix one such fibre $t$.  For each nontrivial divisor $d\mid R$, let $k_d(t)$ be the number of
compatible congruences among the twelve possible moduli
\[
 3^a5^b d,\qquad 0\le a\le3,\quad0\le b\le2.
\]
If we cap every $k_d(t)$ at four, the square-free cofactor classes cover at most
\[
 U_4(R)=\frac{315604}{323323}\approx0.97612604114.
\]
Every fifth or later copy of modulus $d$ can add at most a further fraction $1/d$.  Therefore
coverability of a non-pure fibre forces the weighted excess condition
\begin{equation}\label{eq:heavy-fibre}
 \boxed{
 \sum_{\substack{d\mid R\\d>1}}
 \frac{(k_d(t)-4)_+}{d}
 \ge
 \frac{7719}{323323}
 \approx0.02387395886.}
\end{equation}
Summed over the at least 188 non-pure fibres, the required weighted excess is at least
\[
 \frac{1451172}{323323}\approx4.48830426539.
\]
Thus every hypothetical covering of $L_0$ must manufacture a substantial family of
\emph{heavy base fibres}: fibres on which several of the $3^a5^b d$ classes with the same
square-free cofactor divisor $d$ line up simultaneously.

There is also an immediate finite reduction in the cofactor divisor.  For fixed $d$, at most
12 base multipliers $3^a5^b$ are available, so $(k_d(t)-4)_+\le8$.  The total contribution that
could possibly come from the divisors $d\ge1729$ is
\[
 8\sum_{\substack{d\mid R\\d\ge1729}}\frac1d
 =\frac{6544}{323323}
 <\frac{7719}{323323}.
\]
Hence every non-pure fibre must receive positive heavy excess from one of the following 19
small cofactor divisors:
\[
\begin{split}
&7,11,13,17,19,77,91,119,133,143,187,209,221,247,323,\\
&1001,1309,1463,1547.
\end{split}
\]
More quantitatively, these 19 divisors alone must contribute at least
\[
 \frac{1175}{323323}\approx0.00363413676
\]
of weighted excess on every non-pure fibre.  This cuts the cross-$d$ search from 31 possible
square-free divisors to 19 before any optimization is performed.

This reduction still does not exclude $L_0$.  In fact, an explicit set-valued witness on the
$675$-point base shows that the current heavy-excess relaxation is feasible.  The pure base
classes
\[
0(3),\ 4(5),\ 2(9),\ 13(15),\ 0(25),\ 5(27),\ 8(45),\ 10(75),\
7(135),\ 16(225),\ 1(675)
\]
cover exactly $487$ base fibres, attaining the preceding upper bound and leaving exactly 188.
For each of the 19 small cofactor divisors one can then choose a compatible aligned base residue
profile for the classes $3^a5^b d$ so that, on every one of those 188 uncovered fibres, the
small-divisor weighted excess is at least
\[
 \frac{86}{17017}\approx0.00505376976
 >\frac{1175}{323323}.
\]
The full list of aligned residues is verified exactly in
\texttt{heavy\_set\_relaxation.py}.  This is not a covering of $L_0$; it is a witness that the
current \emph{base-only relaxation} survives.  Consequently the missing obstruction must retain
the residue classes in the square-free coordinate $\Z/R\Z$ as well as their $675$-base support.

A first coarse moment bound does not yet contradict \eqref{eq:heavy-fibre}.  For one fixed
$d$, if $N_d(t)=k_d(t)-1$ counts only the eleven nontrivial base multipliers, then
\[
 (N_d(t)-3)_+\le \frac15\binom{N_d(t)}2.
\]
The sum of the maximum pairwise intersection sizes of the eleven divisor classes on
$\Z/675\Z$ is exactly $525$, hence
\[
 \sum_{t\bmod675}(k_d(t)-4)_+\le105.
\]
After summing this separately over all $d\mid R$, the resulting ceiling is about $52.13$, far
above the required $4.49$.  So the present moment estimate is too weak by roughly an order of
magnitude.  The missing information is not another scalar density: it is the simultaneous
compatibility pattern of the heavy sets for the different square-free divisors $d$.  That is now
the main target of the fibre attack.

The earlier reciprocal-capacity calculation is retained in \texttt{fiber\_relaxation.py} as a
negative diagnostic; all new exact constants and CRT set calculations in this section are
reproduced by \texttt{cofactor\_set\_attack.py}.


\section{Conditioning on the $7$-coordinate of the square-free cofactor}
\label{sec:seven-coordinate}

The $675\times R$ analysis leaves
\[
 R=7S,\qquad
 S=11\cdot13\cdot17\cdot19=46{,}189.
\]
It is useful to absorb the $7$-coordinate into the prime-power base and write
\[
 L_0=AS,
 \qquad
 A=3^3 5^2 7=4{,}725.
\]
This is still a conditioning on the $7$-coordinate of $R$: the point is that the remaining
$S$-coordinate is now a square-free product of four relatively large primes, while all of the
$3$-, $5$-, and $7$-adic branching sits in the much smaller outer group $\Z/A\Z$.

For a necessary-condition argument we may adjoin an arbitrary congruence for every divisor
modulus that is absent from a hypothetical covering.  Adding classes can only help coverage.
Thus partition the full divisor family into three groups:
\begin{enumerate}[label=(\roman*)]
\item pure outer moduli $b\mid A$, $b>1$;
\item pure cofactor moduli $e\mid S$, $e>1$;
\item mixed moduli $be$ with $b>1$ and $e>1$.
\end{enumerate}
Let $O\subseteq\Z/A\Z$ be the complement of the pure outer classes and let
$C\subseteq\Z/S\Z$ be the complement of the pure cofactor classes.

\subsection{A common residual in the four-prime cofactor}

McNew--Setty's CRT inclusion--exclusion lemma applies to a set or multiset of moduli
\cite{McNewSetty2026}.  With one copy of every nontrivial divisor of $S$ it gives
\[
 U_1(S)=\frac{12839}{46189}.
\]
Consequently
\begin{equation}\label{eq:C-lower}
 \boxed{|C|\ge 46189-12839=33350.}
\end{equation}
The important point is set-valued: the same $C$ occurs in \emph{every} one of the $A=4725$
outer fibres because the pure $S$ classes have no $3$-, $5$-, or $7$-component.

The same one-copy calculation on
\[
 A=3^3 5^2 7
\]
gives
\begin{equation}\label{eq:O-lower}
 U_1(A)=\frac{4006}{4725},
 \qquad
 \boxed{|O|\ge719.}
\end{equation}
This is closely related to the finite CRT capacity certificates formalized by
Mian--Siddique \cite{MianSiddique2026}: pairwise-coprime classes force a definite uncovered
set before the remaining divisor budget is charged against it.

\subsection{Residual capacity of a mixed class}

Fix $b\mid A$, $b>1$, and a residue class $B_b\subseteq\Z/A\Z$ modulo $b$.
If a prime $p\in\{3,5,7\}$ does not divide $b$, then the already-present pure class modulo $p$
intersects $B_b$ in exactly a $1/p$ fraction, independently by CRT.  Ignoring every other pure
outer class therefore gives the rigorous upper bound
\begin{equation}\label{eq:alpha-b}
 |B_b\cap O|
 \le
 \alpha(b):=
 \frac{A}{b}
 \prod_{\substack{p\in\{3,5,7\}\\p\nmid b}}
 \left(1-\frac1p\right).
\end{equation}
Likewise, for a residue class $E_e$ modulo a nontrivial divisor $e\mid S$,
\begin{equation}\label{eq:beta-e}
 |E_e\cap C|
 \le
 \beta(e):=
 \frac{S}{e}
 \prod_{\substack{q\in\{11,13,17,19\}\\q\nmid e}}
 \left(1-\frac1q\right)
 =\varphi(S/e).
\end{equation}
The last equality uses square-freeness of $S$.

A mixed class modulo $be$ factors under CRT as $B_b\times E_e$.  Therefore its intersection
with $O\times C$ has size at most $\alpha(b)\beta(e)$.  Exact summation gives
\[
 \sum_{\substack{b\mid A\\b>1}}\alpha(b)=3482,
 \qquad
 \sum_{\substack{e\mid S\\e>1}}\beta(e)=11629.
\]
Hence every mixed class together has total union capacity at most
\begin{equation}\label{eq:mixed-capacity}
 3482\cdot11629=40{,}492{,}178
\end{equation}
points of $O\times C$, even before subtracting any overlap between mixed classes.

Combining \eqref{eq:C-lower} and \eqref{eq:mixed-capacity}, a covering would require
\[
 33350\,|O|\le40{,}492{,}178.
\]
Since
\[
 1214\cdot33350=40{,}486{,}900
 <40{,}492{,}178
 <40{,}520{,}250
 =1215\cdot33350,
\]
we obtain the derived necessary interval
\begin{equation}\label{eq:O-window}
 \boxed{719\le |O|\le1214.}
\end{equation}
Thus any surviving $L_0$ configuration must make its pure $3^3 5^2 7$ subsystem unusually
efficient: it must cover between $3511$ and $4006$ of the $4725$ outer residues.

\subsection{Direct concentration along the $7$-coordinate}

The preceding product argument treats $A$ as one block.  The $7$-coordinate itself gives a
second restriction.  The pure modulus-$7$ class kills one of the seven slices.  On any chosen
$r$ of the other six slices, the $7$-free outer classes contribute one common copy of every
nontrivial divisor of $675$, while the classes with modulus $7d$ contribute one additional copy
of $d$ which may be assigned to one of the chosen slices.

Applying the McNew--Setty multiset bound termwise, and optimizing the copy assignments in the
direction that maximizes the total covered density of the selected slices, gives the following
exact residual totals:
\[
\begin{array}{c|rrrrrr}
 r&1&2&3&4&5&6\\
\hline
 \text{residual lower bound}&-143&-33&155&343&531&719.
\end{array}
\]
The first two values are nonpositive and therefore make no claim.  The last four do.  In
particular, any five of the six nontrivial $7$-slices contain at least $531$ points of $O$.
Consequently
\begin{equation}\label{eq:seven-concentration}
 \boxed{
 \max_{a\bmod7}|O\cap(a\bmod7)|\le |O|-531.}
\end{equation}
For $r=6$ the argument reproduces $|O|\ge719$ from a specifically $7$-adic decomposition.

\subsection{A six-copy bridge through the common $675$-residual}

There is a second way to expose the same $7$-adic structure which is useful for locating the
remaining loss.  Let
\[
 T\subseteq\Z/675\Z
\]
be the common residual left by the $7$-free pure moduli $d\mid675$, $d>1$.  The one-copy CRT
bound used in the preceding section gives
\[
 |T|\ge188,
\]
and the explicit $675$-point witness in Section~\ref{sec:cofactor} shows that this lower
bound is sharp.

For $d\mid675$, $d>1$, write
\[
 M_d(T):=\max_{a\bmod d}|T\cap(a\bmod d)|,
 \qquad
 H(T):=\sum_{\substack{d\mid675\\d>1}}M_d(T).
\]
After the pure modulus-$7$ class removes one $7$-slice, there are six copies of $T$.  Each
local class of modulus $7d$ can delete points from only one of those copies, and it can delete
at most $M_d(T)$ points there.  Hence, by a union bound over the eleven local classes,
\begin{equation}\label{eq:T-lower-O}
 \boxed{|O|\ge6|T|-H(T).}
\end{equation}
Likewise, for a $7$-free probe modulus $d$ the largest intersection with $O$ is at most
$6M_d(T)$, while for a probe modulus $7d$ it is at most $M_d(T)$.  The modulus-$7$ probe
itself contributes at most $|T|$.  Therefore
\begin{equation}\label{eq:T-upper-AO}
 \boxed{A(O)\le |T|+7H(T).}
\end{equation}
Both inequalities are rigorous, but this scalar compression is too coarse to finish the
argument.  An exact fixed residual produced by exploratory local search has
\[
 |T|=212,
 \qquad H(T)=409,
 \qquad \frac{H(T)}{|T|}=1.929245283\ldots.
\]
For that $T$, equations \eqref{eq:T-lower-O}--\eqref{eq:T-upper-AO} give only
$|O|\ge863$ and $A(O)\le3075$.  Thus the next successful step must retain how the eleven
local $7d$ deletions interact with the maximizing residue classes, rather than replacing their
geometry by the two scalars $|T|$ and $H(T)$.  The fixed witness is checked exactly in
\texttt{seven\_coordinate\_attack.py}; no extremality is claimed.

\subsection{The remaining finite lemma on $4725$ points}

For an actual residual $O$ define
\[
 A(O):=
 \sum_{\substack{b\mid4725\\b>1}}
 \max_{a\bmod b}|O\cap(a\bmod b)|,
\]
and, for the common cofactor residual,
\[
 B(C):=
 \sum_{\substack{e\mid46189\\e>1}}
 \max_{c\bmod e}|C\cap(c\bmod e)|.
\]
For fixed $O$ and $C$, every mixed class $be$ contributes at most the product of the
corresponding two maxima, so
\begin{equation}\label{eq:AOBC}
 |O|\,|C|\le A(O)B(C)
\end{equation}
is necessary for the mixed classes to cover $O\times C$.
Equation \eqref{eq:beta-e} gives the universal bound $B(C)\le11629$.  Therefore the single
finite inequality
\begin{equation}\label{eq:outer-target}
 \boxed{
 A(O)<\frac{33350}{11629}|O|
 \approx2.86783042394\,|O|
 }
\end{equation}
for every pure-outer residual $O$ would eliminate $L_0$ outright.  A slightly stronger but
cleaner sufficient target is
\[
 A(O)\le\frac{14}{5}|O|.
\]
This is now the main finite subproblem produced by the $7$-coordinate attack.  It lives on only
$4725$ points and $23$ outer divisor shapes.

\subsection{Exact exploratory witnesses}

The script \texttt{seven\_coordinate\_attack.py} contains two fixed actual-set configurations
found during exploratory local search and verifies them without floating-point decisions.  One
outer assignment has
\[
 |O|=895,
 \qquad
 A(O)=2468,
 \qquad
 \frac{A(O)}{|O|}=2.75754189944<\frac{14}{5}.
\]
A pure-$S$ assignment attains the one-copy CRT coverage bound itself, leaving
\[
 |C|=33350,
 \qquad
 B(C)=11451.
\]
For this concrete pair,
\[
 \frac{A(O)B(C)}{|O||C|}
 =\frac{14130534}{14924125}
 \approx0.94682495624<1.
\]
Thus those two residual geometries cannot be completed by the mixed classes even under the
sum-of-maxima relaxation.  This is exploratory evidence, not a universal exclusion: a
hypothetical covering could choose different pure residue assignments.  Its value is that it
identifies \eqref{eq:outer-target}, rather than a $218$-million-point SAT instance, as the next
concrete proof target.

\section{Relation to the current computational frontier}

Mian and Siddique formalized finite exclusions in Lean and proved, with a kernel-checked
certificate chain, that any odd covering LCM exceeds $10{,}000$ \cite{MianSiddique2026}.  Their
paper emphasizes a finite CRT bridge designed to consume future search certificates.  McNew and
Setty use optimization to push the classification of primitive covering numbers essentially
through $10^6$, with one unresolved even case in that range \cite{McNewSetty2026}.  These works
and the present BFF-based calculation have different trust models and purposes: the aim here is
a compact structural certificate rather than a large black-box classification.

Recent constructive work also shows that the odd-modulus problem is close to solvable when one
permits limited repetition; for example, constructions with only a small number of repeats of a
chosen odd modulus are studied in \cite{Bispels2025}.  Such constructions reinforce the view
that prime-power branching is central rather than incidental.

\section{Next proof targets}

The numerical checkpoint at $L_0$ gives a focused research program.

\begin{enumerate}[label=\arabic*.]
\item \textbf{Outer residual concentration at $A=4725$.}  Prove
\eqref{eq:outer-target}, or the cleaner sufficient inequality
$A(O)\le14|O|/5$, for every residual left by one class for each nontrivial divisor of $4725$.
This would eliminate $L_0$ immediately by the product-residual argument of
Section~\ref{sec:seven-coordinate}.

\item \textbf{Mixed-block overlap.}  Strengthen Lemma~\ref{lem:block} by replacing the union
bound on mixed moduli with an overlap-sensitive bound.  The natural inputs are the CRT
dependency graph and exact block capacities $r(M)$.

\item \textbf{Higher-order BFF clusters.}  Generalize the triangle experiment to a certified
junction-tree or hypergraph inequality using all exactly known intersections of pairwise-disjoint
index sets.  The target is not merely a better numerical correction but a closed-form inequality
that remains monotone in the prime parameters.

\item \textbf{Independent certification.}  Reimplement the 105-case exponent certificate in a
second system and conduct a targeted literature search for the exact $218{,}243{,}025$ corollary
and maximum-spanning-forest formulation before making any novelty claim.
\end{enumerate}

\section{Reproducibility map}

\begin{center}
\footnotesize
\begin{tabularx}{\textwidth}{@{}>{\raggedright\arraybackslash}p{0.36\textwidth}>{\raggedright\arraybackslash}X>{\raggedright\arraybackslash}p{0.19\textwidth}@{}}
\toprule
File & Role & Evidence level\\
\midrule
\texttt{bff\_core.py} & BFF formulas + optimized forest & published + derived\\
\texttt{verify\_218m\_bound.py} & 105 exact exponent cases & exact computation\\
\texttt{certificate\_218m.csv} & human-auditable case table & exact computation\\
\texttt{msf\_refinement.py} & optimized forest edges & derived + exact computation\\
\texttt{triangle\_exploration.py} & one-triangle cluster & derived/exploratory\\
\texttt{tower\_capacity.py} & block-capacity lemma/test & derived/exploratory\\
\texttt{fiber\_relaxation.py} & 3-adic scalar-capacity diagnostic & exploratory/exact\\
\texttt{cofactor\_set\_attack.py} & CRT-aware 3-adic/$675$ cofactor attack & derived + exact\\
\texttt{cofactor\_set\_output.txt} & exact constants from cofactor attack & exact computation\\
\texttt{heavy\_set\_relaxation.py} & explicit 675-base relaxation witness & exploratory + exact\\
\texttt{heavy\_set\_relaxation\_output.txt} & witness residues and checks & exact computation\\
\texttt{seven\_coordinate\_attack.py} & $7$-coordinate/product-residual attack & derived + exact\\
\texttt{seven\_coordinate\_output.txt} & exact constants and fixed witnesses & exact computation\\
\texttt{research\_log.md} & provenance and next steps & audit trail\\
\bottomrule
\end{tabularx}
\end{center}

\bibliographystyle{alpha}
\bibliography{references}

\end{document}
